\section{Adapter-Aware Inference Systems}
\label{sec:inference_systems}

The following section examines how existing systems manage the challenge of
serving multiple task-specific adapters efficiently over a shared frozen
backbone. The literature reviewed here operates exclusively within a single
administrative domain; its relevance to the research direction of this SLR
lies in the serving and memory management primitives it establishes, and in
the centralisation assumptions it makes that a decentralised system would
need to overcome.

\subsection{The Multi-Tenant Serving Challenge}
\label{sec:inference_systems:intro}

Deploying large language models in production requires serving many concurrent users
efficiently. When each user or task is associated with a distinct PEFT adapter, the
serving system must manage a \textit{population} of adapters alongside a single shared
frozen backbone, a problem qualitatively different from single-model serving.
Naively loading a new adapter for each request is prohibitively expensive: a modern
GPU holds only a handful of large model weights in high-bandwidth memory (HBM) at once,
and per-request adapter swaps incur memory transfer latency that dominates over the
inference computation itself for short requests. The multi-tenant adapter serving
literature addresses this challenge through techniques including unified memory
management, batching of heterogeneous-adapter requests, CPU--GPU compute overlap, and
tiered storage hierarchies.

The foundations of efficient LLM serving were established before the adapter era.
\citet{YuOrca2022} introduced iteration-level scheduling with selective batching, which dynamically
combines active requests into GPU batches at each generation step, eliminating the
padding waste of static batching and achieving up to 36.9$\times$ throughput improvement
for single-model generation. \citet{KwonPagedAttention2023} introduced PagedAttention,
which manages the KV-cache using a paged virtual memory abstraction, eliminating
fragmentation and enabling memory-efficient sharing of cached KV states across requests.
PagedAttention achieves 2--4$\times$ higher throughput compared to FasterTransformer and
Orca at equivalent latency targets.
These system primitives underpin most of the adapter-aware serving systems reviewed below.

\subsection{GPU-Resident Multi-Tenant LoRA Serving}
\label{sec:inference_systems:punica}

\citet{Chen2023Punica} introduced Punica, the first system explicitly designed for
multi-tenant LoRA serving on a single GPU. The key contribution is a custom CUDA kernel
(SGMV, Segmented Gather Matrix-Vector) that concatenates LoRA computations for requests with different adapters into
a single GPU operation, decoupling LoRA compute from the base model forward pass.
Because LoRA adapters are defined as rank-$r$ matrix products with $r \ll d$, the
additional per-adapter compute is small enough to be amortised across a batch of
simultaneously active adapters. Punica achieved 12$\times$ higher throughput compared to state-of-the-art LLM serving
systems with only 2ms additional latency per token, establishing the technical
feasibility of multi-tenant LoRA serving on shared GPU infrastructure.

The central limitation of Punica is that all adapters must fit simultaneously in GPU HBM.
For a system serving thousands of tasks, this constraint is immediately prohibitive:
adapter libraries for large-scale deployments require orders of magnitude more memory
than any single GPU can accommodate.

\citet{Sheng2024} addressed the GPU capacity constraint with S-LoRA, a system that
stores adapter weights in a unified CPU--GPU memory pool and dynamically pages adapter
weights between CPU and GPU memory as requests arrive and depart. S-LoRA introduced
two key mechanisms: \textit{Unified Paging}, which maintains a unified memory pool that co-manages both adapter weights of
heterogeneous ranks and KV cache tensors with varying sequence lengths, enabling
fine-grained paged allocation for both simultaneously, and a \textit{stateless scheduler} that selects which adapters to
prefetch into GPU memory based on request queue depth and adapter popularity. S-LoRA
scales to serve over 2,000 distinct LoRA adapters concurrently on a single A100 GPU
with near-zero switching overhead, achieving 1.5--4$\times$ higher throughput than
adapter-unaware baselines at the same GPU utilisation, depending on workload
characteristics \citep{Sheng2024}.

S-LoRA's centralised scheduling policy and unified memory management architecture support
thousands of concurrent adapters but rely on a single administrative domain: all GPUs,
CPU memory, and the adapter catalogue are controlled by one provider. This centralisation
is fundamentally incompatible with P2P infrastructure, where adapter discovery involves
a distributed lookup, and where adapter weights may be spread across autonomous peer
nodes with internet-scale latency. This limitation identifies the need for DHT-based
discovery and a P2P exchange framework.

Among centralised serving systems, S-LoRA's architecture is the closest structural analogue to what a decentralised setting would require. A conceptual parallel exists between its Unified Paging hierarchy (GPU HBM - CPU DRAM) and a P2P setting where the secondary tier would be remote peer storage rather than local CPU memory. Whether this analogy holds under real network latency — where bandwidth is variable and prefetching decisions must account for peer churn — is an open question that S-LoRA's design does not address.

\subsection{CPU-Assisted and Rank-Aware Serving}
\label{sec:inference_systems:cpu}

\citet{Li2024CaraServe} proposed CaraServe, which augments S-LoRA by offloading LoRA
computations to CPU cores rather than merely using the CPU as slow memory. Because LoRA
adapter matrices are small (a rank-4 adapter for a hidden dimension 4096 occupies only
$2 \times 4 \times 4096 \times 4 \approx 256$~KB per layer), their matrix products can
be computed efficiently on CPU while the GPU processes the base model forward pass.
CaraServe early-starts adapter prefilling on CPU while adapter weights are concurrently
transferred to GPU, switching to GPU execution once loading completes, improving
average request serving latency by up to 1.4$\times$ and achieving up to 99\% SLO
attainment. The rank-aware placement policy: high-rank adapters remain on GPU,
low-rank adapters are computed on CPU, provides a template for hardware-aware adapter
placement in heterogeneous P2P nodes, where different peers may have different compute
capabilities.

\citet{LiToppings2025} proposed Toppings, a system that identifies that under continuous
batching, loading a newly requested adapter to GPU interrupts the ongoing decoding of all
inflight queries. Toppings resolves this via a pipeline loading scheme that simultaneously
coordinates CPU prefilling, GPU transfer, and ongoing decoding, improving average request
latency by up to 1.7$\times$ and achieving 99\% SLO attainment.

\subsection{Serverless and Tiered Storage Serving}
\label{sec:inference_systems:serverless}

Moving beyond in-process memory management, several systems extended the storage
hierarchy to include SSD and object storage. \citet{ChenFlashServe2025} proposed
FlashServe, a serverless LLM inference framework that tiers adapter weights across GPU
HBM, CPU DRAM, and NVMe SSD, using a predictive prefetching policy trained on
historical request patterns to minimise cache miss latency. \citet{FuServerlessLLM2024}
addressed cold-start latency in serverless LLM deployments, introducing fast model
checkpoint loading from distributed storage that achieves sub-second model initialisation.
Both systems demonstrate the feasibility of on-demand adapter loading from persistent
storage under real serving latency budgets.

\citet{NiPredictive2025} introduced a predictive adapter prefetching system that models
request arrival patterns using time-series forecasting to decide which adapters to
preload into GPU memory ahead of their expected use, further reducing cache miss
penalties under skewed request distributions.

\subsection{Edge Deployment and Constrained Hardware}
\label{sec:inference_systems:edge}

A growing body of work addresses adapter serving on resource-constrained edge hardware.
\citet{Shen2025EdgeLoRA} proposed EdgeLoRA, a multi-tenant LoRA serving system for edge
devices (Jetson, Raspberry Pi) that combines adapter compression (quantisation and
structured pruning) with dynamic time-sharing of GPU memory across multiple tenant
adapters. EdgeLoRA achieves up to 70\% memory savings compared to uncompressed loading
while maintaining within-2\% accuracy on representative NLP tasks.

\citet{CaiEdgeLLM2024} introduced a collaborative edge--cloud LLM inference framework
in which edge devices handle lightweight adapter-augmented forward passes for latency-
sensitive requests while offloading full-model computation to cloud servers for
accuracy-critical queries. \citet{ZhangEdgeShard2025} extended edge deployment to
pipeline-parallel sharding: the transformer layers are distributed across a chain of
edge devices, each holding a subset of layers and their associated adapters.

\subsection{Distributed and Collaborative Inference}
\label{sec:inference_systems:distributed}

\citet{Borzunov2023} introduced Petals, a collaborative inference system for large
language models in which geographically distributed volunteers contribute GPU
resources to host contiguous subsets (shards) of transformer layers, with possible overlap for redundancy. Users run the
embedding and final projection layers locally and route activations through the
volunteer network for mid-stack processing. Petals demonstrated that sub-cloud hardware
can collectively serve models that no single participant could host, achieving
throughput sufficient for interactive generation despite commodity network latencies.
Petals is the closest prior work to the P2P inference direction examined in this review,
but it does not incorporate per-task adapter modules: all participating nodes serve
identical backbone layers without task specialisation.

\citet{HuDistributed2022} addressed the bandwidth cost of split inference, proposing
adapter-based bottleneck compression for split-computing architectures in which only
a compressed representation is transmitted between device tiers rather than the full
intermediate activation tensor.

\subsection{Advanced Multi-Adapter Orchestration}
\label{sec:inference_systems:advanced}

Several recent systems have addressed further aspects of large-scale adapter serving.
\citet{WuRock2025} proposed Rock, a multi-adapter serving system for multimodal LLMs
that coordinates thousands of domain-specific adapters across a GPU cluster using a
centralised orchestration layer. \citet{ZhuCannikin2025} introduced Cannikin, which
eliminates straggler effects in multi-LoRA batched inference by decoupling the prefill
and decode phases and routing them to separate GPU resources. \citet{Shi2025AuLoRA}
proposed AuLoRA, which extends unified paging to fine-grained layer-level adapter
loading, enabling sub-adapter granularity in memory management.

Across these systems, a consistent architectural pattern emerges: the adapter weight is
treated as the atom of resource management, and system optimisations operate at the
granularity of individual adapter loading, computation, and eviction events. This
design principle transfers directly to a P2P exchange architecture, where the adapter is also
the atom of network transfer, storage, and composition.

\subsection{Positioning against the Decentralised Setting}
\label{sec:inference_systems:gap}

The existing serving literature addresses the multi-tenant inference problem with considerable technical depth — but it leaves one fundamental question unanswered: where does the adapter catalogue come from, and who controls it? Every system reviewed in this chapter inherits its adapter collection from a centralised upstream process that is never designed, only assumed.
Translating this model to a P2P setting exposes four open problems for which the serving literature provides no tools: how to discover which peers hold relevant adapters without a central registry; how to establish trust in adapters from unknown, autonomous peers — and whether such trust is meaningful at all; how to fetch adapters under real network latency conditions where bandwidth is not guaranteed; and how to incentivise peers to contribute compute and storage without central governance.

The Unified Paging concept introduced by S-LoRA is, in principle, transferable to a P2P setting — remote peer storage replaces CPU DRAM as the secondary memory tier. But that transfer only becomes meaningful once the four open problems above have been resolved.
These four open problems define the systems-level gap analysed in Section~\ref{sec:synthesis_gap}.
